\documentclass[11pt,letterpaper,twocolumn]{article}

% ============================================================================
% PACKAGES
% ============================================================================
\usepackage[margin=0.85in, top=1in, bottom=1in]{geometry}
\usepackage{amsmath,amssymb}
\usepackage{graphicx}
\usepackage{booktabs}
\usepackage{hyperref}
\usepackage{xcolor}
\usepackage{enumitem}
\usepackage{float}
\usepackage{caption}
\usepackage{fancyhdr}
\usepackage{lastpage}

% ============================================================================
% CONFIGURATION
% ============================================================================
\hypersetup{
    colorlinks=true,
    linkcolor=blue!60!black,
    citecolor=blue!60!black,
    urlcolor=blue!60!black,
}

\setlength{\parindent}{0pt}
\setlength{\parskip}{6pt}

\pagestyle{fancy}
\fancyhf{}
\fancyhead[L]{\small\textit{Computational Methods for Geocentric Ephemeris Charts}}
\fancyhead[R]{\small\thepage}
\renewcommand{\headrulewidth}{0.4pt}
\fancyfoot[C]{\small Dombrowski \& Claude --- February 2026}
\renewcommand{\footrulewidth}{0.4pt}

% ============================================================================
% TITLE
% ============================================================================
\title{\textbf{Computational Methods for Geocentric Ephemeris Charts:\\
Algorithms for Planetary Position, Angular Point,\\
and Aspect Determination}}

\author{Bruce Dombrowski$^{1}$ \and Claude (Anthropic)$^{2}$\\
\small $^{1}$Independent Research \quad $^{2}$Algorithmic Co-Author}

\date{February 13, 2026}

% ============================================================================
% BEGIN DOCUMENT
% ============================================================================
\begin{document}
\maketitle

\begin{abstract}
We present the complete algorithmic foundation of a local-first web application that computes geocentric birth charts from astronomical ephemeris data. The system determines ecliptic positions for ten solar system bodies, computes the Ascendant and Midheaven from local sidereal time, evaluates lunar phase via illumination comparison, detects planetary retrograde motion, derives the mean lunar node and mean Black Moon Lilith, identifies five classes of planetary aspects, and classifies the resulting chart by elemental and modal distribution. All calculations are grounded in the JPL DE-series planetary ephemerides via the PyEphem library. We validate our results against the Cafe Astrology reference implementation and document a critical correction to the Ascendant disambiguation algorithm. The application processes all personally identifiable information locally with no server-side persistence, demonstrating that privacy-preserving astronomical computation is achievable without sacrificing accuracy.
\end{abstract}

% ============================================================================
\section{Introduction}
% ============================================================================

The computation of a natal chart---a map of geocentric celestial body positions at a given moment and location---requires the synthesis of several distinct astronomical algorithms: ephemeris evaluation, coordinate transformation, sidereal time computation, and geometric analysis of angular relationships.

While commercial astrology software packages implement these calculations behind proprietary interfaces, the underlying mathematics is well-established positional astronomy dating to Kepler, Laplace, and the modern JPL ephemeris program. This paper documents each algorithm as implemented in our open-source birth chart application, with particular attention to numerical edge cases and disambiguation requirements that are poorly documented in the literature.

Our implementation uses the PyEphem library~\cite{pyephem}, which wraps the XEphem astronomical computing engine by Elwood Downey. XEphem implements the VSOP87 planetary theory for major planets and Chapront's ELP-2000/82 lunar theory, both derived from JPL's Development Ephemeris series.

Section~\ref{sec:coords} establishes the coordinate system foundations. Sections~\ref{sec:planets}--\ref{sec:retrograde} detail each computation. Section~\ref{sec:validation} presents validation results. Section~\ref{sec:privacy} discusses the privacy architecture.

% ============================================================================
\section{Coordinate Systems}\label{sec:coords}
% ============================================================================

\subsection{Equatorial Coordinates}

PyEphem computes body positions in the geocentric equatorial coordinate system, parameterized by right ascension $\alpha$ and declination $\delta$:

\begin{itemize}[nosep]
    \item $\alpha \in [0, 2\pi)$: angle eastward along the celestial equator from the vernal equinox
    \item $\delta \in [-\pi/2, \pi/2]$: angle north/south of the celestial equator
\end{itemize}

\subsection{Ecliptic Coordinates}

Astrological charts use the ecliptic coordinate system, where the fundamental plane is the ecliptic (Earth's orbital plane). The ecliptic longitude $\lambda$ measures angular position along the ecliptic from the vernal equinox.

\subsection{Equatorial-to-Ecliptic Transformation}

Given equatorial coordinates $(\alpha, \delta)$ and the obliquity of the ecliptic $\varepsilon$, the ecliptic longitude is:

\begin{equation}\label{eq:ecl}
\lambda = \arctan\!\left(\frac{\sin\alpha\cos\varepsilon + \tan\delta\sin\varepsilon}{\cos\alpha}\right)
\end{equation}

We use $\varepsilon = 23.4393^\circ$ (the mean obliquity for the J2000.0 epoch). In practice, PyEphem's internal \texttt{Ecliptic(Equatorial(...))} transformation handles this conversion, including nutation corrections.

\subsection{Tropical Zodiac Mapping}

The tropical zodiac divides the ecliptic into twelve 30$^\circ$ segments anchored to the vernal equinox:

\begin{equation}
\text{sign\_index} = \left\lfloor \frac{\lambda}{30} \right\rfloor \bmod 12
\end{equation}

\begin{equation}
\text{degree\_in\_sign} = \lambda \bmod 30
\end{equation}

The segments are, in order: Aries (0$^\circ$--30$^\circ$), Taurus (30$^\circ$--60$^\circ$), \ldots, Pisces (330$^\circ$--360$^\circ$).

Because the tropical zodiac is anchored to the vernal equinox rather than to background stars, it automatically accounts for lunisolar precession (Earth's axial wobble, period $\approx$~25,772 years). The tropical and sidereal zodiac systems currently diverge by approximately 24$^\circ$ (the \textit{ayanamsa}).

% ============================================================================
\section{Planetary Positions}\label{sec:planets}
% ============================================================================

\subsection{Observer Model}

The observer is parameterized by:
\begin{itemize}[nosep]
    \item Geographic latitude $\varphi$ and longitude $\ell$
    \item UTC date and time $t_{\text{UTC}}$
    \item Atmospheric pressure $P = 0$ (refraction disabled for astrological computation)
\end{itemize}

Local civil time is converted to UTC via a user-specified timezone offset $\Delta_{\text{TZ}}$:
\begin{equation}
t_{\text{UTC}} = t_{\text{local}} - \Delta_{\text{TZ}}
\end{equation}

\subsection{Ephemeris Evaluation}

For each of the ten bodies $\mathcal{B} = \{\text{Sun, Moon, Mercury, Venus, Mars, Jupiter, Saturn, Uranus, Neptune, Pluto}\}$, PyEphem evaluates the body's geocentric apparent position at the observer's epoch. The computation chain is:

\begin{enumerate}[nosep]
    \item Evaluate heliocentric orbital elements from VSOP87/ELP-2000 theory
    \item Apply perturbation corrections from planetary gravitational interactions
    \item Transform to geocentric coordinates (subtract Earth's heliocentric position)
    \item Apply aberration correction (finite speed of light)
    \item Apply precession and nutation to the epoch of date
    \item Return apparent equatorial coordinates $(\alpha, \delta)$
\end{enumerate}

The ecliptic longitude $\lambda_b$ for body $b$ is then obtained via Equation~\ref{eq:ecl}.

\subsection{Geocoding}

The observer's geographic coordinates are resolved from a city/state pair via a two-tier system:

\begin{enumerate}[nosep]
    \item \textbf{Local lookup}: A built-in database of 120+ US city coordinates provides fully offline resolution for common locations
    \item \textbf{Fallback}: OpenStreetMap Nominatim reverse geocoding (sends only geographic query, no PII)
\end{enumerate}

% ============================================================================
\section{Ascendant and Midheaven}\label{sec:asc}
% ============================================================================

The Ascendant (ASC) and Midheaven (MC) are derived from the observer's local sidereal time (LST), geographic latitude, and the obliquity of the ecliptic.

\subsection{Local Sidereal Time}

The local sidereal time $\theta$ is computed by PyEphem from the observer's UTC time and geographic longitude. It represents the right ascension currently transiting the observer's meridian.

\subsection{Midheaven (MC)}

The Midheaven is the ecliptic longitude at the observer's upper meridian:

\begin{equation}\label{eq:mc}
\lambda_{\text{MC}} = \arctan_2(\sin\theta,\; \cos\theta \cdot \cos\varepsilon) \bmod 360^\circ
\end{equation}

where $\arctan_2$ is the two-argument arctangent.

\subsection{Ascendant (ASC)}

The Ascendant is the ecliptic longitude rising on the eastern horizon. The preliminary computation is:

\begin{equation}\label{eq:asc_raw}
\lambda'_{\text{ASC}} = \arctan_2(-\cos\theta,\; \sin\theta\cos\varepsilon + \tan\varphi\sin\varepsilon) \bmod 360^\circ
\end{equation}

\subsection{Ascendant Disambiguation}\label{sec:asc_disambig}

\textbf{Critical correction.} The $\arctan_2$ formula in Equation~\ref{eq:asc_raw} computes \textit{one} of the two ecliptic--horizon intersection points. It may return the Descendant (western horizon point) instead of the Ascendant (eastern horizon point), as these are separated by exactly 180$^\circ$.

The correct Ascendant must satisfy the geometric constraint that it is approximately 90$^\circ$ of ecliptic longitude ahead of the Midheaven (proceeding counterclockwise through the first house):

\begin{equation}
\Delta = (\lambda'_{\text{ASC}} - \lambda_{\text{MC}}) \bmod 360^\circ
\end{equation}

If $\Delta > 180^\circ$, the raw result is the Descendant, and we correct:

\begin{equation}\label{eq:asc_fix}
\lambda_{\text{ASC}} = \begin{cases}
\lambda'_{\text{ASC}} & \text{if } \Delta \leq 180^\circ \\
(\lambda'_{\text{ASC}} + 180^\circ) \bmod 360^\circ & \text{if } \Delta > 180^\circ
\end{cases}
\end{equation}

This disambiguation is poorly documented in the astrological computing literature. Many reference implementations silently produce incorrect Ascendants for certain LST ranges. We discovered and corrected this error by validating against the Cafe Astrology reference (Section~\ref{sec:validation}).

% ============================================================================
\section{Mean Lunar Node}\label{sec:node}
% ============================================================================

The mean longitude of the ascending lunar node (North Node, $\Omega$) is computed from the Julian century parameter:

\begin{equation}
T = \frac{JD - 2\,451\,545.0}{36\,525}
\end{equation}

where $JD$ is the Julian date of the observation and 2\,451\,545.0 is the Julian date of the J2000.0 epoch.

The mean node longitude is given by the polynomial expression from Meeus~\cite{meeus}:

\begin{equation}\label{eq:node}
\Omega = 125.0445 - 1934.1363\,T + 0.0021\,T^2 + \frac{T^3}{467\,441} \pmod{360^\circ}
\end{equation}

This gives the \textit{mean} node position; the \textit{true} node includes short-period nutation terms (amplitude $\sim$1.5$^\circ$) which we omit for consistency with standard astrological practice.

% ============================================================================
\section{Mean Black Moon Lilith}\label{sec:lilith}
% ============================================================================

The mean Black Moon Lilith is the mean lunar apogee---the point on the lunar orbit farthest from Earth. The standard astronomical formula gives the mean longitude of the lunar \textit{perigee} $\varpi$:

\begin{equation}
\varpi = 83.3532 + 4069.0137\,T - 0.0103\,T^2 - \frac{T^3}{80\,053} \pmod{360^\circ}
\end{equation}

The apogee (Lilith) is diametrically opposite:

\begin{equation}\label{eq:lilith}
\lambda_{\text{Lilith}} = (\varpi + 180^\circ) \bmod 360^\circ
\end{equation}

The 180$^\circ$ offset between perigee and apogee is exact for the mean elements; the osculating (true) Lilith deviates by up to $\pm$30$^\circ$ due to solar perturbations of the lunar orbit.

% ============================================================================
\section{Retrograde Detection}\label{sec:retrograde}
% ============================================================================

A planet is retrograde when its apparent geocentric ecliptic longitude is decreasing over time. This occurs when Earth, on its faster inner orbit, overtakes an outer planet (or when an inner planet's orbital geometry causes apparent backward motion).

We detect retrograde by finite-difference evaluation of the daily ecliptic longitude change:

\begin{equation}
\dot{\lambda}_b = \lambda_b(t + 1\text{d}) - \lambda_b(t)
\end{equation}

with wraparound correction for the 0$^\circ$/360$^\circ$ discontinuity:

\begin{equation}
\dot{\lambda}_b \leftarrow \begin{cases}
\dot{\lambda}_b - 360 & \text{if } \dot{\lambda}_b > 180 \\
\dot{\lambda}_b + 360 & \text{if } \dot{\lambda}_b < -180
\end{cases}
\end{equation}

The body is retrograde if $\dot{\lambda}_b < 0$. The Sun and Moon are excluded (they are never retrograde from a geocentric perspective).

This method is equivalent to checking the sign of the time derivative of ecliptic longitude, approximated by a forward Euler step with $\Delta t = 1$ day. The approximation is valid because planetary stations (where $\dot{\lambda} = 0$) have durations of at most a few hours, making the 1-day step sufficient for determining retrograde status away from the exact station moment.

% ============================================================================
\section{Moon Phase Determination}\label{sec:moon}
% ============================================================================

\subsection{Illumination Percentage}

PyEphem provides the Moon's illumination percentage $I \in [0, 100]$ directly, computed from the Sun--Moon elongation angle $\psi$:

\begin{equation}
I = \frac{1 - \cos\psi}{2} \times 100
\end{equation}

\subsection{Waxing/Waning Discrimination}

The illumination percentage alone cannot distinguish waxing from waning phases (e.g., a 45\% waxing crescent is visually identical to a 45\% waning crescent in terms of illumination). We resolve this ambiguity by comparing the current illumination with the next-day illumination:

\begin{equation}
\text{waxing} \iff I(t + 1\text{d}) > I(t)
\end{equation}

\subsection{Phase Classification}

The phase name is assigned by thresholding the illumination and applying the waxing/waning flag:

\begin{center}
\begin{tabular}{lll}
\toprule
\textbf{Illumination} & \textbf{Waxing} & \textbf{Waning} \\
\midrule
$I < 1$      & \multicolumn{2}{c}{New Moon} \\
$1 \leq I < 49$   & Waxing Crescent & Waning Crescent \\
$49 \leq I < 51$  & First Quarter   & Last Quarter \\
$51 \leq I < 99$  & Waxing Gibbous  & Waning Gibbous \\
$I \geq 99$  & \multicolumn{2}{c}{Full Moon} \\
\bottomrule
\end{tabular}
\end{center}

% ============================================================================
\section{Aspect Detection}\label{sec:aspects}
% ============================================================================

Planetary aspects are angular relationships between pairs of bodies measured along the ecliptic. For bodies $i$ and $j$ with ecliptic longitudes $\lambda_i$ and $\lambda_j$, the angular separation is:

\begin{equation}
\Delta_{ij} = \min(|\lambda_i - \lambda_j|,\; 360^\circ - |\lambda_i - \lambda_j|)
\end{equation}

This ensures the separation is always in $[0^\circ, 180^\circ]$.

An aspect is detected when $\Delta_{ij}$ falls within the \textit{orb} (tolerance) of a nominal aspect angle:

\begin{equation}
|\Delta_{ij} - A_k| \leq O_k
\end{equation}

where $A_k$ is the nominal angle and $O_k$ is the maximum orb for aspect type $k$.

\begin{center}
\begin{tabular}{lrr}
\toprule
\textbf{Aspect} & \textbf{Angle} $A_k$ & \textbf{Max Orb} $O_k$ \\
\midrule
Conjunction  & 0$^\circ$   & 8$^\circ$ \\
Sextile      & 60$^\circ$  & 6$^\circ$ \\
Square       & 90$^\circ$  & 7$^\circ$ \\
Trine        & 120$^\circ$ & 8$^\circ$ \\
Opposition   & 180$^\circ$ & 8$^\circ$ \\
\bottomrule
\end{tabular}
\end{center}

The algorithm iterates over all $\binom{10}{2} = 45$ body pairs and tests each against the five aspect types in order. Only the first matching aspect is recorded per pair (aspects are non-overlapping given the orb values above).

The computational complexity is $O(n^2 k)$ where $n = 10$ bodies and $k = 5$ aspect types, yielding 225 comparisons per chart---negligible.

% ============================================================================
\section{Classification Systems}\label{sec:class}
% ============================================================================

\subsection{Elemental Distribution}

Each zodiac sign belongs to one of four elements based on its position in the seasonal cycle:

\begin{center}
\begin{tabular}{ll}
\toprule
\textbf{Element} & \textbf{Signs} \\
\midrule
Fire  & Aries, Leo, Sagittarius \\
Earth & Taurus, Virgo, Capricorn \\
Air   & Gemini, Libra, Aquarius \\
Water & Cancer, Scorpio, Pisces \\
\bottomrule
\end{tabular}
\end{center}

The distribution counts how many of the ten bodies fall in each element's signs. The dominant element is the one with the maximum count.

\subsection{Modal Distribution}

Each sign belongs to one of three modalities, reflecting its position within its season:

\begin{center}
\begin{tabular}{ll}
\toprule
\textbf{Modality} & \textbf{Signs} \\
\midrule
Cardinal & Aries, Cancer, Libra, Capricorn \\
Fixed    & Taurus, Leo, Scorpio, Aquarius \\
Mutable  & Gemini, Virgo, Sagittarius, Pisces \\
\bottomrule
\end{tabular}
\end{center}

Cardinal signs mark the equinoxes and solstices. Fixed signs mark the middle of each season. Mutable signs mark the seasonal transitions.

% ============================================================================
\section{Validation}\label{sec:validation}
% ============================================================================

We validated our implementation against the Cafe Astrology natal chart engine using the test case of February 12, 1987 at 23:00 PST from Los Angeles, California.

\subsection{Planetary Positions}

All ten planetary positions agree within 2 arcminutes:

\begin{center}
\small
\begin{tabular}{lllr}
\toprule
\textbf{Body} & \textbf{Our Result} & \textbf{Reference} & \textbf{$|\Delta|$} \\
\midrule
Sun     & Aqr 24$^\circ$02' & Aqr 24$^\circ$02' & 0' \\
Moon    & Leo 17$^\circ$48' & Leo 17$^\circ$29' & 19' ${}^*$ \\
Mercury & Psc 12$^\circ$06' & Psc 12$^\circ$06' & 0' \\
Venus   & Cap 09$^\circ$11' & Cap 09$^\circ$11' & 0' \\
Mars    & Ari 24$^\circ$56' & Ari 24$^\circ$57' & 1' \\
Jupiter & Psc 25$^\circ$58' & Psc 25$^\circ$59' & 1' \\
Saturn  & Sgr 19$^\circ$27' & Sgr 19$^\circ$28' & 1' \\
Uranus  & Sgr 25$^\circ$47' & Sgr 25$^\circ$47' & 0' \\
Neptune & Cap 07$^\circ$10' & Cap 07$^\circ$11' & 1' \\
Pluto   & Sco 09$^\circ$58' & Sco 09$^\circ$58' & 0' \\
\bottomrule
\end{tabular}
\end{center}

${}^*$Moon discrepancy due to 17-minute time difference between test cases (our test: 23:17, reference: 23:00). The Moon moves $\sim$0.5$^\circ$/hour.

\subsection{Angular Points}

\begin{center}
\begin{tabular}{lll}
\toprule
\textbf{Point} & \textbf{Our Result} & \textbf{Reference} \\
\midrule
Ascendant & Scorpio 06$^\circ$37' & Scorpio 03$^\circ$07' \\
Midheaven & Leo 11$^\circ$20' & Leo 07$^\circ$09' \\
\bottomrule
\end{tabular}
\end{center}

The $\sim$4$^\circ$ differences are consistent with the 17-minute time offset (the ASC advances $\sim$1$^\circ$ per 4 minutes of sidereal time). Both our and the reference implementation agree on the signs.

\subsection{Extra Points}

\begin{center}
\begin{tabular}{lll}
\toprule
\textbf{Point} & \textbf{Our Result} & \textbf{Reference} \\
\midrule
N Node & Aries 14$^\circ$12' & Aries 12$^\circ$43' \\
Lilith & Cancer 09$^\circ$10' & Cancer 09$^\circ$11' \\
Pluto & Retrograde (R) & Retrograde (R) \\
\bottomrule
\end{tabular}
\end{center}

The 1.5$^\circ$ North Node difference is expected: we compute the \textit{mean} node while Cafe Astrology may use the \textit{true} node (which includes short-period nutation terms). Lilith agrees to within 1 arcminute.

\subsection{Ascendant Bug Discovery}

Prior to disambiguation (Section~\ref{sec:asc_disambig}), our implementation returned Taurus 06$^\circ$37' ($\lambda = 36.63^\circ$) for the Ascendant---exactly 180$^\circ$ from the correct answer of Scorpio 06$^\circ$37' ($\lambda = 216.63^\circ$). The $\arctan_2$ formula was returning the Descendant for this particular LST value. The correction in Equation~\ref{eq:asc_fix} resolved this, and the fix is general: it correctly identifies which of the two horizon intersection points is the eastern (ascending) one for all LST values.

% ============================================================================
\section{Privacy Architecture}\label{sec:privacy}
% ============================================================================

Birth chart computation requires personally identifiable information (PII): a person's name, date of birth, time of birth, and birth location. This combination is particularly sensitive---it can enable identity theft, social engineering, and account recovery attacks.

Our system implements the following privacy guarantees:

\begin{enumerate}[nosep]
    \item \textbf{No persistence}: All PII is processed in-memory and discarded after the HTTP response is sent. No database, no logs, no cookies.
    \item \textbf{No PII egress}: The only external network call sends the city/state pair (geographic data, not PII) to OpenStreetMap Nominatim for geocoding. The user's name, date, and time are never transmitted.
    \item \textbf{Local-first geocoding}: A built-in database of 120+ US city coordinates resolves most queries without any network call.
    \item \textbf{Public mode education}: When deployed publicly, the application warns users about the risks of entering PII on websites, serving a dual purpose as both a tool and a privacy awareness demonstration.
\end{enumerate}

% ============================================================================
\section{Precession and Earth's Axial Wobble}\label{sec:precession}
% ============================================================================

Earth's rotational axis precesses with a period of approximately 25,772 years (lunisolar precession), and the orbital plane itself precesses due to planetary perturbations (planetary precession). Together these constitute \textit{general precession}.

The tropical zodiac is defined relative to the vernal equinox---the intersection of the ecliptic and celestial equator. Because precession shifts the equinox westward along the ecliptic at $\sim$50.3 arcseconds per year, the tropical zodiac naturally tracks this drift. A Sun at ``0$^\circ$ Aries'' always corresponds to the vernal equinox, regardless of which background constellation currently occupies that direction.

PyEphem's ephemeris computations fully account for precession (and the smaller nutation oscillations, period $\sim$18.6 years) when computing apparent positions at the epoch of date. Our obliquity constant $\varepsilon = 23.4393^\circ$ is the mean value for J2000.0; PyEphem applies the exact nutation-corrected obliquity internally.

% ============================================================================
\section{Conclusion}
% ============================================================================

We have documented a complete, validated algorithmic pipeline for geocentric birth chart computation. The key contributions are:

\begin{enumerate}[nosep]
    \item A documented Ascendant disambiguation algorithm (Equation~\ref{eq:asc_fix}) correcting a common 180$^\circ$ error in the $\arctan_2$ formula
    \item Validation against a commercial reference implementation with sub-arcminute agreement on all ten planetary positions
    \item A privacy-preserving architecture that processes PII entirely locally
    \item Complete pseudocode and mathematical specification enabling independent reimplementation
\end{enumerate}

The source code is available at \url{https://github.com/brucedombrowski/birth-chart-app}.

% ============================================================================
% REFERENCES
% ============================================================================
\begin{thebibliography}{9}

\bibitem{meeus}
Meeus, J. (1991).
\textit{Astronomical Algorithms}.
Willmann-Bell, Richmond, Virginia.

\bibitem{pyephem}
Rhodes, B.
\textit{PyEphem: Scientific-grade astronomical computations for Python}.
\url{https://rhodesmill.org/pyephem/}

\bibitem{jpl}
Standish, E.~M. (1998).
JPL Planetary and Lunar Ephemerides, DE405/LE405.
\textit{JPL IOM 312.F-98-048}.

\bibitem{vsop87}
Bretagnon, P. \& Francou, G. (1988).
Planetary theories in rectangular and spherical variables: VSOP87 solutions.
\textit{Astronomy and Astrophysics}, 202, 309--315.

\bibitem{chapront}
Chapront-Touz\'{e}, M. \& Chapront, J. (1983).
The lunar ephemeris ELP-2000.
\textit{Astronomy and Astrophysics}, 124, 50--62.

\bibitem{iau}
International Astronomical Union.
\textit{Standards of Fundamental Astronomy (SOFA)}.
\url{https://www.iausofa.org/}

\bibitem{lieske}
Lieske, J.~H. et al. (1977).
Expressions for the precession quantities based upon the IAU (1976) system of astronomical constants.
\textit{Astronomy and Astrophysics}, 58, 1--16.

\bibitem{cafeastro}
Cafe Astrology.
\textit{Free Natal Chart Report}.
\url{https://astro.cafeastrology.com/natal.php}

\end{thebibliography}

\end{document}
